% HD-S48. GPT-6 Astra Pro /root/sixfolds, inherited effort; 21 November 2026.
\section{HD-S48: the monomial rank-three normal subgroup}
\label{sec:monomial-rank-three-scope}

Let \(K\subset\mathbf C\) have characteristic zero, and let
\(N\subset GL_m(K)\), \(m\ge15\), be finite and absolutely
irreducible. Suppose that \(N\) is generated by actual
moving-rank-three elements, contains no nonidentity element of
moving rank one or two, and has a complex system of lines.
Write \(P\) for its permutation action on those lines.
When \(P\) is primitive, its diagonal kernel \(D\) is
characteristic and nonscalar, and all the original line projectors
belong to \(M_m(K)\). Every actual \(K\)-linear overgroup
normalizing \(N\) preserves this descended line system.

In the sixfold setting, take
\[
 K=\mathbf Q(i),\quad\hbox{or}\quad
 K=\mathbf Q(\sqrt{-d}),\qquad
 d\in\{2,7,11,19,43,67,163\},
 \tag{S48.1}
\]
as the exact character field of the original finite coefficient
group \(G\subset GL_m(K)\), and suppose \(N\triangleleft G\).
Use its natural primitive sixfold of signature \((3,3)\),
from a connected full-monodromy sphere cover, with the
restricted Jacobian theta polarization.
For primitive \(P\), the seven latter fields are impossible;
for \(K=\mathbf Q(i)\), the original natural primitive
sixfold has ordinary Hermitian determinant class \([-1]\).
Together with HD-S46 and HD-S47, this removes the entire
complex-imprimitive absolutely irreducible \(N_3\) alternative
from the remaining determinant sources of HD-R15.
The structural assertion does not assume that the lines were
defined over \(K\), nor that a scalar central extension splits.

\paragraph{1. The actual permutation image has small support.}
Choose nonzero vectors in the complex lines. On a cycle of
length \(\ell\), a monomial matrix has fixed dimension one
when the product of its cycle entries is one, and zero otherwise:
the fixed-vector equations successively determine all coordinates
from the first and impose exactly this final return equation.
Thus a nontrivial cycle moves rank at least \(\ell-1\).
For a moving-rank-three generator, its nontrivial permutation
cycles therefore satisfy
\[
 \sum_{\ell\ge2}(\ell-1)\le3,
 \qquad
 \sum_{\ell\ge2}\ell\le
          2\sum_{\ell\ge2}(\ell-1)\le6.
 \tag{S48.2}
\]
The action on the line set is transitive, since the sum over
any orbit would otherwise be a proper invariant subspace.
Some generator has nonidentity permutation, because the action
is transitive in degree greater than one. Hence the minimum
nonidentity support size of \(P\), traditionally its
\emph{class}, is at most six.

We use two precise published class bounds. Manning,
\emph{The order of primitive groups. III}, Transactions AMS19
(1918), 127--142, section~38, printed p.142, states that primitive
groups of class four and six have degree at most eight and ten,
respectively. This is the sentence immediately before
TheoremXV, not a specialization of the numerical hypothesis
of that theorem:
\url{https://doi.org/10.1090/S0002-9947-1918-1501093-5}.
Jordan, \emph{Recherches sur les substitutions}, Journal de
math\'ematiques pures et appliqu\'ees, series2,17 (1872),
351--367, section~24, printed p.367, classifies the primitive groups
of prime class \(q>3\): for \(q\) not of the form
\(2^a-1\), the sole group of class \(q\) is cyclic of
degree \(q\). We use only \(q=5\):
\url{https://www.numdam.org/item/JMPA_1872_2_17__351_0.pdf}.
Both displayed original pages were read directly. These
published classification inputs are not being relabeled as
new classification proofs.

For completeness, the classes two and three have an elementary
argument. All conjugate transpositions generate the symmetric
groups on the connected components of their edge graph.
These components are blocks for a primitive group, so the graph
is connected and the group contains \(S_m\).
For three-cycles use the hypergraph of their supports.
The subgroup generated by all three-cycles is normal and
nontrivial, so is transitive in a primitive group; its support
hypergraph is connected. Three-cycles on a connected hypergraph
generate the full alternating group on its vertices. Here is
an explicit induction: two supports intersecting in two vertices
extend the already generated alternating group by the third
vertex. If the new support meets the old vertex set in just
\(a\), choose distinct old vertices \(b,c\); for
\(x=(a\,b\,c)\) and \(y=(a\,d\,e)\), the actual
commutator is
\[
                    xyx^{-1}y^{-1}=(a\,b\,d).
 \tag{S48.3}
\]
This first adds \(d\) by the two-vertex case and then adds
\(e\) using \(y\). The two-vertex case follows by
conjugating \((a\,b\,d)\) by the alternating group on
the old vertices and using its inverse; these give the
three-cycles with the new vertex and any two old vertices.
Connectivity completes the induction.

Class one is impossible for a nonidentity permutation. The
published bounds exclude classes four, five and six in degree
\(m\ge15\). Consequently, for primitive \(P\),
\[
                         P=A_m\quad\hbox{or}\quad P=S_m.
 \tag{S48.4}
\]

\paragraph{2. A scalar diagonal kernel gives an actual bireflection.}
Suppose \(D\) were scalar. Let \(H_i\subset N\)
stabilize the first specified line \(L_i\), and let
\(\chi_i:H_i\to\mathbf C^*\) be its actual action there.
The character is faithful on \(D\), because a scalar acting
as one on a nonzero line is the identity. It kills \(H_i'\).
The permutation projection \(\pi:N\to P\) is onto and
\(\pi(H_i)=P_i\), so
\[
 H_i'\cap D=1,\qquad
 \pi(H_i')=P_i'=A_{m-1},\qquad
          \pi|_{H_i'}:H_i'\xrightarrow{\sim}A_{m-1}.
 \tag{S48.5}
\]
Choose the actual element \(g\in H_i'\) projecting to a
three-cycle on labels other than \(i\). It has order three
by this injective map. For every other label \(j\) fixed
by that cycle, the stabilizer of \(L_j\) in \(H_i'\)
is exactly \(A_{m-2}\). It is perfect; hence its actual
one-dimensional action on \(L_j\) is trivial. The action
on \(L_i\) is also trivial, since \(\chi_i\) kills
\(H_i'\). On the three moved lines, \(g\) is one
weighted three-cycle and \(g^3=1\) makes its return scalar
one. The return equation in paragraph1 now gives
\[
                            \operatorname{rank}(g-I)=2.
 \tag{S48.6}
\]
This is an actual forbidden element of \(N\), not a
statement merely about its projective image. Thus \(D\)
is nonscalar. This argument includes all scalar central
extensions of (S48.4).

\paragraph{3. Characteristicness and exact Galois descent.}
Every abelian normal subgroup of \(N\) has abelian normal
image in \(A_m\) or \(S_m\). Neither has a nontrivial
abelian normal subgroup for \(m\ge5\): simplicity and
perfection of \(A_m\), and its trivial centralizer in
\(S_m\), give this assertion. The image is therefore
trivial, so every abelian normal subgroup lies in \(D\).
Since \(D\) itself is abelian and normal, it is the unique
largest such subgroup and is characteristic in \(N\).

Let \(\chi_j:D\to\mathbf C^*\) be its characters on
the specified lines. Equality of two characters defines an
equivalence relation on the line set, invariant under \(P\).
Primitivity implies that all characters are equal or that
they are pairwise distinct. Equality of all would make every
element of \(D\) scalar, which was excluded. Thus they are
distinct. Put \(E=K(\zeta_{\exp D})\). The literal
projectors are
\[
 p_j=\frac1{|D|}\sum_{d\in D}\chi_j(d^{-1})\rho(d),
 \qquad p_j^2=p_j,\quad p_jp_k=0\ (j\ne k),
 \quad\sum_jp_j=I,\quad\operatorname{rank}p_j=1.
 \tag{S48.7}
\]
Character orthogonality proves every identity and identifies
their images with the original complex lines. They are also
the complete set of primitive central idempotents of
\(E[D]\) in its given matrix image.

The finite Galois group \(\operatorname{Gal}(E/K)\)
permutes these idempotents. Its permutation commutes with the
actual conjugation action of every matrix in \(N\), whose
entries belong to \(K\); hence it commutes with (S48.4).
The centralizer of \(A_m\) in \(S_m\) is trivial:
a commuting permutation must send label \(i\) to a label
fixed by its point stabilizer \(A_{m-1}\), whose only
fixed label is \(i\). Therefore every Galois element fixes
each \(p_j\), and
\[
                p_j\in M_m(K),\qquad
                U_j=p_jK^m,\qquad\dim_KU_j=1.
 \tag{S48.8}
\]
Because \(D\) is characteristic, every actual normalizer
of \(N\) in \(GL_m(K)\) permutes the same primitive
central projectors and their \(K\)-lines. This proves the
structural assertion with explicit projectors, without
assuming Galois stability of a chosen block system.

\paragraph{4. The original coefficient module and cyclic quotient.}
Return to the exact-field sixfold setting (S48.1), and let
\(H=\operatorname{Stab}_G(U_1)\). Choose
\(0\ne v_1\in U_1\), \(t_jU_1=U_j\), and
\(v_j=\rho(t_j)v_1\). The line character
\(\chi:H\to K^*\) is finite, and the actual matrix rule is
\[
 \rho(g)v_j=
       \chi(t_{g(j)}^{-1}gt_j)v_{g(j)}.
 \tag{S48.9}
\]
The map \(t_j\otimes v_1\mapsto v_j\) is the actual
isomorphism \(\operatorname{Ind}_H^G\chi\simeq V\).
In pushforward-function form, with the literal projector
\(p_1\), its inverse pair is
\[
 \begin{aligned}
 F_v(x)&=p_1\rho(x)v,&F_v(hx)&=\chi(h)F_v(x),\\
 F_{\rho(g)v}(x)&=F_v(xg),&
 F&\longmapsto\sum_{H\backslash G}\rho(x)^{-1}F(x).
 \end{aligned}
 \tag{S48.10}
\]
The sheet label \(Ht_j^{-1}\leftrightarrow t_jH\)
relates this right monodromy to the original left deck action.
On an inertia cycle a fixed section is exactly a vector fixed
by its return character, so these maps extend to the invariant
stalks and the parabolic cohomology of the ramified quotient.

For the seven fields with \(d\in\{2,7,11,19,43,67,163\}\),
the finite character values in (S48.9) are \(\pm1\).
The formula gives rational trace for every \(g\in G\),
contrary to its exact character field. This excludes those
cases through the actual induced matrices.

For \(K=\mathbf Q(i)\), the values lie in \(\mu_4\),
and \(\chi(H)=\mu_4\); a smaller image would again give
rational trace. Choose \(h_0\in H\) with \(\chi(h_0)=i\),
put \(H_0=\ker\chi\), and form the actual curves
\[
  Y=\widetilde C/H,\qquad X=\widetilde C/H_0,
       \qquad\alpha:\widetilde C\longrightarrow X.
 \tag{S48.11}
\]
The map \(X\to Y\) is cyclic of degree four. In the
unit monomial basis (S48.9) let \(h(v_j,v_k)=\delta_{jk}\),
take \(v=(1+i)v_1\) of norm two, and set \(D_0=|H|=|G|/m\).
The rational corner projector is literally
\[
 \begin{aligned}
 e_v&=\frac m{|G|}\sum_{g\in G}
       \operatorname{Re}h(v,\rho(g^{-1})v)g\\
    &=\frac1{|H|}\sum_{h_1\in H}
           (\chi(h_1)^{-1}+\bar\chi(h_1)^{-1})h_1
      =\operatorname{av}_{H_0}\frac{1-h_0^2}{2}.
 \end{aligned}
 \tag{S48.12}
\]
Off-\(H\) coefficients vanish. Absolute irreducibility of
the induced module makes this a primitive rational corner with
field \(K\), as follows also from matrix-coefficient
orthogonality. The integer multiple \(D_0e_v\) is an actual
endomorphism of \(J\widetilde C\). With covariant left-deck
actions \(a_g=L_{g,*}\), its connected image \(B_v\)
has integral Gaussian action \(i_{B_v}=a_{h_0}^{-1}|_{B_v}\).
Indeed \(h_0\) commutes with (S48.12) and
\(e_vh_0^2=-e_v\).

If \(B_X\) is the cyclic primitive factor of \(JX\),
(S48.12) makes \(\alpha^*|_{B_X}\) an actual surjective
isogeny onto \(B_v\), with
\[
 \alpha_*\alpha^*=[|H_0|],\qquad
       (\alpha^*)^*\theta_{J\widetilde C}=|H_0|\theta_{JX}.
 \tag{S48.13}
\]
Thus the quotient realizes the same primitive Hodge factor and
field action, with a specified polarization multiplier.
For the real coefficient map
\(\Phi=\operatorname{PD}(j+\bar j)\),
\(j_x(g)=h(v,\rho(g)x)/2\), the complex Schur factor is
\(|G|/(2m)\) and the two real cross terms give exactly
\[
                         \Phi^*h_{B_v}=D_0 H_V.
 \tag{S48.14}
\]
This retains the actual normalization from HD-S23/HD-S34.

The original prescribed primitive line may be \(Kw\),
with \(w\notin U_1\). Set \(q_w=h(w,w)\) and
\(A(x)=v h(w,x)/q_w\). Its actual rational Fourier operator
is
\[
 b=\frac{2m}{|G|}\sum_{g\in G}
           \operatorname{Re}\operatorname{tr}
                       (A\rho(g^{-1}))g,
       \qquad b^\dagger b=\frac2{q_w}e_w.
 \tag{S48.15}
\]
Matrix-coefficient orthogonality gives \(A\), its conjugate,
and zero on the remaining coefficient summands. Clearing
denominators by \(M\) therefore gives an actual isogeny
\(B_w\to B_v\) with pullback polarization multiplier
\(2M^2/q_w\). Any originally specified positive invariant
coefficient form is a positive rational multiple of \(h\)
by Schur's lemma; this scalar cancels in the corner and the
ratio in (S48.15). Every positive rational multiplier changes
a rank-six Hermitian determinant by its sixth power, the norm
of its cube. Thus these maps retain the ordinary determinant
class of the originally specified factor.

HD-S34's integral parabolic argument now applies to the
literal matrices (S48.9), without any restriction on local
cycle lengths. Over \(R=\mathbf Z[i,1/2]\) their lattice
is self-dual. A local cycle has holonomy \(\zeta\in\mu_4\)
and coinvariants \(R/(\zeta-1)\); the nontrivial holonomies
give zero, and holonomy one has the explicit invariant and
coinvariant generators with perfect pairing. Global
coinvariants vanish through the actual relation
\((i-1)[v_1]=0\) from \(h_0\).
The integral parabolic pairing is therefore perfect by the
absolute/relative cap-product maps of HD-S34. Its rank-six
Hermitian determinant is a negative rational unit of \(R\),
namely \(-2^a\). Since \(2=N(1+i)\), (S48.14)--(S48.15)
give class \([-1]\) for the original factor, with the split
isometry supplied by HD-S03a. Compatible flat polarization
changes in a rank-nine family retain the class by HD-S17's
explicit common-centralizer calculation; no arbitrary isolated
polarization change is asserted.

\paragraph{5. Exhaustion of the complex-imprimitive alternative.}
If \(P\) is imprimitive, group the original lines over any
nontrivial permutation block system. This gives an actual
complex block system for \(N\), each block being the direct
sum of the corresponding lines. Let its block dimension be
\(b>1\). Transitivity ensures that some rank-three generator
permutes these blocks; its nontrivial cycle has moving rank at
least \(b\) by the same return-map calculation. Hence
\(b\le3\). For \(b=2\), \(m\ge15\) gives
\(m=2r\) with \(r\ge8\), exactly the hypotheses of
HD-S47, including its actual characteristic kernel and
polarized quotient maps; its period rank is at most seven.
For \(b=3\), \(r\ge5\), and HD-S46 descends the
original three-block system through its characteristic kernel.
Its explicit quotient/corner maps retain the original factor;
period rank nine forces a genus-zero rank-three coefficient
system generated by twelve reflections, whose determinant
source is already covered by HD-S31/HD-S40 and HD-G27.
For primitive \(P\), paragraphs1--4 apply. These alternatives
exhaust every complex monomial \(N\). Together with the
already proved two- and three-dimensional block cases, they
close the entire complex-imprimitive form3 of HD-R15.

\paragraph{6. Verification scope.}
The author read the complete proof, including the actual scalar
kernel lift, the Galois projectors, the induction functions,
the ramified quotient, and the original-line Fourier isogeny.
The only new published class inputs are the exact original
Manning and Jordan statements identified in paragraph1.
No finite numerical computation is required for this universal
argument, and none is being presented as a classification.
